\documentclass[11pt]{article}

\usepackage[margin=0.75in,headheight=14pt]{geometry}
\usepackage[T1]{fontenc}
\usepackage{lmodern}
\usepackage{microtype}
\usepackage[table]{xcolor}
\usepackage{array}
\usepackage{amssymb}
\usepackage{tabularx}
\usepackage{longtable}
\usepackage{enumitem}
\usepackage{titlesec}
\usepackage{fancyhdr}
\usepackage{lastpage}

\definecolor{Primary}{HTML}{1F4E79}
\definecolor{Secondary}{HTML}{EAF2F8}
\definecolor{Accent}{HTML}{D9EAF7}
\definecolor{RuleGray}{HTML}{D9D9D9}
\definecolor{TextGray}{HTML}{3A3A3A}

\pagestyle{fancy}
\fancyhf{}
\lhead{Minimum Curriculum Quality Checklist}
\rhead{Curriculum Review Handout}
\cfoot{Page \thepage\ of \pageref{LastPage}}
\renewcommand{\headrulewidth}{0.4pt}
\renewcommand{\footrulewidth}{0pt}

\titleformat{\section}{\large\bfseries\color{Primary}}{}{0pt}{}
\titleformat{\subsection}{\normalsize\bfseries\color{TextGray}}{}{0pt}{}
\setlist[itemize]{leftmargin=1.1em, itemsep=2pt, topsep=2pt}
\setlist[enumerate]{leftmargin=1.4em, itemsep=3pt, topsep=3pt}
\renewcommand{\arraystretch}{1.16}
\setlength{\parindent}{0pt}
\setlength{\parskip}{4pt}
\setlength{\LTpre}{4pt}
\setlength{\LTpost}{8pt}
\emergencystretch=2em

\newcolumntype{Y}{>{\raggedright\arraybackslash}X}
\newcolumntype{L}[1]{>{\raggedright\arraybackslash}p{#1}}
\newcolumntype{C}[1]{>{\centering\arraybackslash}p{#1}}
\newcommand{\checkbox}{\(\square\)}

\begin{document}

{\Huge\bfseries\color{Primary} Minimum Curriculum Quality Checklist}\par
\vspace{3pt}
{\large A practical review handout for evaluating whether a curriculum meets a baseline standard for quality, equity, engagement, and continuous improvement.}\par
\vspace{8pt}
\color{RuleGray}\hrule\color{black}
\vspace{10pt}

\rowcolors{2}{Secondary}{white}
\begin{tabularx}{\textwidth}{>{\bfseries\color{Primary}}p{1.45in}Y}
Purpose & Use this checklist to determine whether a curriculum is effective, student-centered, inclusive, logically sequenced, and adaptable over time. \\
Review Focus & Learning outcomes, content structure, equity, relevance, engagement, assessment, teacher usability, and revision practices. \\
How to Use & Mark each item Yes, No, or Needs Work. A curriculum should meet all critical acceptance criteria before it is considered ready for implementation or continued use. \\
\end{tabularx}
\rowcolors{2}{}{}

\section{Minimum Checklist}

\footnotesize
\rowcolors{2}{Secondary}{white}
\setlength{\tabcolsep}{2pt}
\begin{longtable}{L{0.25in} L{1.50in} L{3.10in} C{0.45in} C{0.45in} C{0.58in}}
\rowcolor{Primary}
\textcolor{white}{\#} & \textcolor{white}{Checklist Item} & \textcolor{white}{Minimum Standard / Evidence to Verify} & \textcolor{white}{Yes} & \textcolor{white}{No} & \textcolor{white}{Needs Work} \\
\endfirsthead
\rowcolor{Primary}
\textcolor{white}{\#} & \textcolor{white}{Checklist Item} & \textcolor{white}{Minimum Standard / Evidence to Verify} & \textcolor{white}{Yes} & \textcolor{white}{No} & \textcolor{white}{Needs Work} \\
\endhead
1 & Clear learning outcomes & The curriculum defines specific, measurable outcomes for each unit, module, course, or major learning sequence. & \checkbox & \checkbox & \checkbox \\
2 & Logical scope and sequence & Content progresses from foundational knowledge to more advanced concepts, skills, applications, and performance expectations. & \checkbox & \checkbox & \checkbox \\
3 & Student-centered design & Lessons, activities, pacing, and supports are designed around student readiness, learning needs, developmental levels, and goals. & \checkbox & \checkbox & \checkbox \\
4 & Inclusive and equitable access & Materials account for diverse cultures, backgrounds, abilities, language needs, and access constraints so all students can participate meaningfully. & \checkbox & \checkbox & \checkbox \\
5 & Real-world relevance & Content connects to students' lives, communities, future careers, civic participation, or authentic social and technical problems. & \checkbox & \checkbox & \checkbox \\
6 & Engagement and active learning & Students have structured opportunities to discuss, investigate, create, collaborate, solve problems, and apply knowledge. & \checkbox & \checkbox & \checkbox \\
7 & 21st-century skills & The curriculum develops critical thinking, creativity, communication, collaboration, digital literacy, and responsible use of technology. & \checkbox & \checkbox & \checkbox \\
8 & Flexible and adaptable structure & The curriculum can be adjusted based on student performance, teacher feedback, technology changes, community needs, and new evidence. & \checkbox & \checkbox & \checkbox \\
9 & Balanced assessment strategy & Assessment includes formative checks, performance tasks, projects, reflections, discussions, and summative measures aligned to outcomes. & \checkbox & \checkbox & \checkbox \\
10 & Student reflection & Students have regular opportunities to reflect on what they learned, how they learned it, and what they need to improve. & \checkbox & \checkbox & \checkbox \\
11 & Collaborative development & Teachers, subject matter experts, instructional leaders, and relevant stakeholders contribute to curriculum design, validation, or review. & \checkbox & \checkbox & \checkbox \\
12 & Continuous improvement process & There is a defined cycle for reviewing student evidence, gathering feedback, revising materials, and documenting improvements. & \checkbox & \checkbox & \checkbox \\
13 & Representation and perspective & The curriculum serves as both a mirror and a window: students see themselves represented and encounter perspectives different from their own. & \checkbox & \checkbox & \checkbox \\
14 & Instruction-assessment alignment & Learning outcomes, instructional activities, practice opportunities, and assessments are visibly connected. & \checkbox & \checkbox & \checkbox \\
15 & Teacher usability & Teachers can understand, implement, adapt, and explain the curriculum without excessive ambiguity or hidden assumptions. & \checkbox & \checkbox & \checkbox \\
\end{longtable}
\setlength{\tabcolsep}{6pt}
\rowcolors{2}{}{}
\normalsize

\section{Minimum Acceptance Criteria}

A curriculum should not be considered ready unless the following critical criteria are satisfied.

\begin{enumerate}
  \item Learning outcomes are clear, observable, and measurable.
  \item Content follows a logical sequence that supports cumulative learning.
  \item The curriculum is inclusive, accessible, and responsive to diverse student needs.
  \item Students actively engage with material through discussion, collaboration, inquiry, creation, and application.
  \item Assessment is ongoing, balanced, and aligned to stated outcomes.
  \item The curriculum develops critical thinking, collaboration, communication, creativity, and digital literacy.
  \item There is a documented process for revision, stakeholder review, and continuous improvement.
\end{enumerate}

\section{Review Rating Guide}

\rowcolors{2}{Secondary}{white}
\begin{tabularx}{\textwidth}{>{\bfseries\color{Primary}}p{2.05in}Y}
\rowcolor{Accent}
Rating & Meaning \\
Meets Minimum Standard & All critical acceptance criteria are met, and checklist gaps are minor, documented, and manageable. \\
Needs Improvement & One or more important checklist items are incomplete, unclear, weakly evidenced, or inconsistently implemented. \\
Not Acceptable & Learning outcomes, assessment alignment, equity, sequencing, or teacher usability are missing or materially weak. \\
\end{tabularx}
\rowcolors{2}{}{}

\section{Reviewer Notes}

\begin{tabularx}{\textwidth}{p{1.8in}Y}
Curriculum / Course: & \vspace{0.28in} \\
Reviewer: & \vspace{0.28in} \\
Review Date: & \vspace{0.28in} \\
Overall Rating: & \checkbox\ Meets Minimum Standard \quad \checkbox\ Needs Improvement \quad \checkbox\ Not Acceptable \\
Priority Improvements: & \vspace{0.75in} \\
Decision / Next Step: & \vspace{0.55in} \\
\end{tabularx}

\end{document}



